% !TEX TS-program = xelatex
% !TEX encoding = UTF-8 Unicode
% !Mode:: "TeX:UTF-8"

\documentclass{resume}
\usepackage{zh_CN-Adobefonts_external} % Simplified Chinese Support using external fonts (./fonts/zh_CN-Adobe/)
% \usepackage{NotoSansSC_external}
% \usepackage{NotoSerifCJKsc_external}
% \usepackage{zh_CN-Adobefonts_internal} % Simplified Chinese Support using system fonts
\usepackage{linespacing_fix} % disable extra space before next section
\usepackage{cite}

\begin{document}
\pagenumbering{gobble} % suppress displaying page number

\name{你的大名}

\basicInfo{
  \email{yuanbin2014@gmail.com} \textperiodcentered\ 
  \phone{(+86) 131-221-87xxx} \textperiodcentered\ 
  \linkedin[billryan8]{https://www.linkedin.com/in/billryan8}}
 
\section{\faGraduationCap\  教育背景}
\datedsubsection{\textbf{电子科技大学}，四川成都}{2023 -- 至今}
\textit{在读硕士研究生}\ 信息与通信工程, 预计 2026 年 6 月毕业
\datedsubsection{\textbf{电子科技大学}, 四川成都}{2019 -- 2023}
\textit{学士}\ 信息与软件工程

\section{\faUsers\ 实习/项目经历}
\datedsubsection{\textbf{华为技术有限公司} 四川成都}{2022年1月 -- 2022年6月}
\role{软件开发工程师}{实习}
\begin{onehalfspacing}
实习期间主要完成了以下工作：
\begin{itemize}
  \item Trace 日志开发
  \item DT FUZZ 测试
  \item 算法题讲解
  \item 自动部署脚本
\end{itemize}
\end{onehalfspacing}

\datedsubsection{\textbf{gcs}}{2024年7月 -- 至今}
\role{Java，TypeScript}{团队项目}
\begin{onehalfspacing}
后端仓库：https://github.com/CMIPT/gcs-back-end。

前端仓库：https://github.com/CMIPT/gcs-front-end。

该项目为自建的 Git 中心仓库，主要用于托管教研室的代码。
项目由本人发起并主导，主要负责项目的设计以及开发。
项目目前在使用中，目前支持以下功能：
\begin{itemize}
  \item 用户注册、登录
  \item SSH 配置
  \item 公有、私有仓库
  \item 仓库合作者
  \item 仓库不同 Ref 预览
\end{itemize}
\end{onehalfspacing}

\datedsubsection{\textbf{GTNDSystem}}{2024年10月 -- 2025年5月}
\role{Python}{团队项目}
\begin{onehalfspacing}
项目原使用 Python2 和 Django 进行开发，主要用于网络拓扑的探测及可视化。
本人负责项目的重构，使用 Python3 按照 PEP8 规范进行重构，同时定义相关开发相关规范。
\end{onehalfspacing}

\datedsubsection{\textbf{blink-cmp-git}}{2025年1月 -- 至今}
\role{Lua}{个人项目}
\begin{onehalfspacing}
项目仓库：https://github.com/Kaiser-Yang/blink-cmp-git。

blink-cmp-git 是一个 Neovim 的插件，该插件提供 Git 中心仓库相关的补全
（Issues、Pull Requests、Commits 以及 Mentions）。
\end{onehalfspacing}

\datedsubsection{\textbf{blink-cmp-dictionary}}{2025年1月 -- 至今}
\role{Lua}{个人项目}
\begin{onehalfspacing}
项目仓库：https://github.com/Kaiser-Yang/blink-cmp-dictionary。

blink-cmp-dictionary 是一个 Neovim 的插件，该插件提供了单词补全。
\end{onehalfspacing}

\datedsubsection{\textbf{blink-cmp-avante}}{2025年1月 -- 至今}
\role{Lua}{个人项目}
\begin{onehalfspacing}
项目仓库：https://github.com/Kaiser-Yang/blink-cmp-avante。

blink-cmp-avante 提供 Avante 插件中内置命令的补全。
\end{onehalfspacing}

\datedsubsection{\textbf{win-resizer.nvim}}{2025年1月 -- 至今}
\role{Lua}{个人项目}
\begin{onehalfspacing}
项目仓库：https://github.com/Kaiser-Yang/blink-cmp-win-resizer.nvim。

win-resizer.nvim 对 Neovim 的 resize 命令进行了封装，提供了一些简单易用的 API。
\end{onehalfspacing}

\datedsubsection{\textbf{t-SNE}}{2024年3月 -- 2024年8月}
\role{C++，Python，Bash}{个人项目}
\begin{onehalfspacing}
项目仓库：https://github.com/Kaiser-Yang/t-SNE。

使用 C++ 对 t-SNE 算法进行实现，使用 Python 对分类结果进行可视化。
\end{onehalfspacing}

\datedsubsection{\textbf{matrix-calculation-accelerator}}{2024年3月 -- 2024年7月}
\role{C++}{个人项目}
\begin{onehalfspacing}
项目仓库：https://github.com/Kaiser-Yang/matrix-calculation-accelerator。

一个矩阵计算加速器，使用 C++ 多线程进行加速。
\end{onehalfspacing}

\datedsubsection{\textbf{vni}}{2023年10月 -- 2024年4月}
\role{C}{个人项目}
\begin{onehalfspacing}
项目仓库：https://github.com/Kaiser-Yang/vni。

一个 Linux 的内核模块。可以统计到指定地址的 IP 包的数量。
\end{onehalfspacing}

% Reference Test
%\datedsubsection{\textbf{Paper Title\cite{zaharia2012resilient}}}{May. 2015}
%An xxx optimized for xxx\cite{verma2015large}
%\begin{itemize}
%  \item main contribution
%\end{itemize}

\section{\faCogs\ IT 技能}
% increase linespacing [parsep=0.5ex]
\begin{itemize}[parsep=0.5ex]
  \item 编程语言：C/C++ > Java > Python == Bash > JS/TS > Dart
  \item 平台：Linux / Windows + WSL
  \item 工具：Neovim / VSCode
\end{itemize}

\section{\faHeartO\ 获奖情况}
\datedline{\textit{二等奖}, 2025华为软件精英挑战赛西南赛区复赛}{2025年4月}
\datedline{\textit{第一名}, 2025华为软件精英挑战赛西南赛区初赛}{2025年3月}
\datedline{\textit{全球前5\%}，IEEEXtreme 18.0 Programming Competition}{2025年3月}
\datedline{\textit{优秀研究生}，电子科技大学}{2024年11月}
\datedline{\textit{二等奖}，电子科技大学2024-2025学年研究生学业奖学金}{2024年11月}
\datedline{\textit{三等奖}, 第九届电子科技大学IEEEXtreme极限编程邀请赛}{2024年7月}
\datedline{\textit{三等奖}，第十三届蓝桥杯全国软件和信息技术专业人才大赛全国总决赛}{2022年5月}
\datedline{\textit{一等奖}，第十三届蓝桥杯全国软件和信息技术专业人才大赛四川赛区}{2022年5月}
\datedline{\textit{华储星}，华为技术有限公司存储部门}{2022年4月}
\datedline{\textit{最佳代码奖}，华为技术有限公司存储部门}{2022年4月}
\datedline{\textit{标兵奖学金}，电子科技大学2021-2022年度}{2022年11月}
\datedline{\textit{标兵奖学金}，电子科技大学2020-2021年度}{2021年11月}
\datedline{\textit{标兵奖学金}，电子科技大学2019-2020年度}{2020年11月}
\datedline{\textit{一等奖}，NOIP2018提高组四川赛区}{2018年12月}

\section{\faInfo\ 其他}
% increase linespacing [parsep=0.5ex]
\begin{itemize}[parsep=0.5ex]
  \item 博客：https://kaiser-yang.github.io 和 https://blog.csdn.net/qq\_45523675。
  \item GitHub：https://github.com/Kaiser-Yang。
\end{itemize}

%% Reference
%\newpage
%\bibliographystyle{IEEETran}
%\bibliography{mycite}
\end{document}
